\section{Method}
\label{sec:method}

We instantiate the decomposition from \S\ref{sec:decomp} in two spaces:
the residual stream, where it defines a block-level observable, and the
attention value space, where it isolates an attention-side source of parallel
updates. These two views lead to the component-scaling interventions used in
the experiments. Figure~\ref{fig:overview} summarizes the resulting analysis
pipeline and how it connects to editing, compression, and pretraining.

\begin{figure*}[t]
  \centering
  \includegraphics[width=0.99\textwidth]{figs/overview/parallel_updates_main_figure_analysis_v11_midline_arrows.pdf}
  \caption{\textbf{Overview of the method and its downstream use.}
           We first measure residual-update geometry at the block level, then
           analyze attention aggregation in value space before output
           projection. This yields a unified component-scaling form for value-
           and residual-space edits, which we use in the experiments on
           inference-time editing, compression error, and pretraining
           behavior.}
  \label{fig:overview}
\end{figure*}

% ------------------------------------------------------------------
\subsection{Residual-Level Decomposition}
\label{sec:residual-level}

\paragraph{Instantiation.}
Set $\Delta = \Du{l}{t}$ (the residual update of block $f^l$ at token $t$)
and $\mathbf{z} = \hz{l}{t}$ (the pre-update hidden state).
For the bounded intervention and localization diagnostic, we report
calibration-corpus averages of the per-token fraction of update magnitude in
the parallel component:
\begin{equation}
  \begin{aligned}
  \pr{l}{t}
    &= r\!\left(\Du{l}{t},\,\hz{l}{t}\right)
     = \frac{|\alpha^l_t|\,\|\hz{l}{t}\|}{\|\Du{l}{t}\|},\\
  \alpha^l_t
    &= \frac{\Du{l}{t}\!\cdot\!\hz{l}{t}}{\|\hz{l}{t}\|^2}.
  \end{aligned}
  \label{eq:para-ratio}
\end{equation}
This residual-level view serves as the shared observable for the rest of the
paper. Figure~\ref{fig:profiles} instead displays the complementary unbounded
ratio \(\|\Delta_{\parallel}\|/\|\Delta_{\perp}\|\); the two quantities are
not numerically interchangeable. The resulting diagnostic is block-agnostic.
We apply it to
(i)~the $\Attn$ sub-layer output,
(ii)~the $\FFN$ sub-layer output, and
(iii)~their joint effect at the transformer-block level.

% ------------------------------------------------------------------
\subsection{Value-Space Decomposition}
\label{sec:value-level}

\paragraph{Instantiation.}
Inside attention, each source token $s$ is mapped to a value vector
$\mathbf{v}_s=W_V\hz{}{s}$. For query token $t$, attention aggregates
these value vectors before the output projection:
\begin{equation}
  \mathbf{o}_t=\sum_{s\le t} A_{ts}\mathbf{v}_s .
  \label{eq:value-aggregation}
\end{equation}
Here $\mathbf{o}_t$ denotes the pre-$W_O$ attention output in value space.
We decompose this aggregate relative to the token's own value direction:
\begin{equation}
  \mathbf{o}_{t,\parallel}
    = \frac{\mathbf{o}_t^\top\mathbf{v}_t}{\|\mathbf{v}_t\|^2}\mathbf{v}_t,
  \qquad
  \mathbf{o}_{t,\perp}=\mathbf{o}_t-\mathbf{o}_{t,\parallel}.
  \label{eq:value-space-decomp}
\end{equation}
This value-space view measures how much of the pre-projection attention message
follows the self-value direction versus other value-space directions.
Residual-level geometry alone cannot identify this structure because the
output projection has already mixed the aggregated value message into the
residual stream.

\paragraph{Raw and exclude-self variants.}
Equation~\ref{eq:value-space-decomp} defines the \emph{raw} variant, which
decomposes the complete value aggregate. Our primary \emph{exclude-self}
variant instead separates the direct self message from the non-self context:
\begin{equation}
  \begin{aligned}
  \mathbf{o}_t
    &= \underbrace{A_{tt}\mathbf{v}_t}_{\mathbf{d}_t\;\text{(self)}}
     + \underbrace{\sum_{s<t}A_{ts}\mathbf{v}_s}_{\mathbf{c}_t\;\text{(non-self)}},\\
  \mathbf{c}_{t,\parallel}
    &= \frac{\mathbf{c}_t^\top\mathbf{v}_t}{\|\mathbf{v}_t\|^2}\mathbf{v}_t,
  \qquad
  \mathbf{c}_{t,\perp}=\mathbf{c}_t-\mathbf{c}_{t,\parallel}.
  \end{aligned}
  \label{eq:exclude-self-decomp}
\end{equation}
We scale only this non-self decomposition and then restore the unchanged self
message:
\begin{equation}
  \widetilde{\mathbf{o}}^{\mathrm{excl}}_t
    = \mathbf{d}_t
      +s^{(\parallel)}\mathbf{c}_{t,\parallel}
      +s^{(\perp)}\mathbf{c}_{t,\perp}.
  \label{eq:exclude-self-scaling}
\end{equation}
Thus, ``exclude-self'' names the aggregate being decomposed: it neither masks
the self-attention edge nor renormalizes the attention row. For multi-head
attention, this split and decomposition are applied independently within each
head before $W_O$.

We apply this value-space analysis specifically to attention. The feed-forward
block also produces residual-level parallel and perpendicular components, but
its intermediate activation has a different channel dimension and is not an
attention aggregation space. We therefore treat MLP outputs at the residual
level, while using value-space decomposition for attention-side updates.

% ------------------------------------------------------------------
\subsection{Component Scaling and Diagonal Form}
\label{sec:closed-form-control}

For the raw value-space and residual-space interventions at query position
$t$, we express component suppression or amplification with a general
parallel-perpendicular rescaling form:
\begin{equation}
  \begin{aligned}
  \mathbf{y}_t
    &= \sum_{s\le t} A_{ts}\mathbf{m}_s,\\
  \mathbf{y}_{t,\parallel}
    &= \frac{\mathbf{y}_t^\top\mathbf{r}_t}{\|\mathbf{r}_t\|^2}\mathbf{r}_t,
  \qquad
  \mathbf{y}_{t,\perp}=\mathbf{y}_t-\mathbf{y}_{t,\parallel}.
  \end{aligned}
  \label{eq:unified-component-decomp}
\end{equation}
where $\mathbf{m}_s$ is the message vector being aggregated and
$\mathbf{r}_t$ is the reference vector whose direction is used to define parallelness. The
value- and residual-space interventions use different choices of these
quantities
($\Delta_t^{\Attn}$ denotes the post-$W_O$ attention branch output):
\begin{equation}
  \begin{aligned}
    \text{V-space:}\quad
    &\mathbf{m}_s=\mathbf{v}_s,\quad
      \mathbf{r}_t=\mathbf{v}_t,\\
    &\mathbf{y}_t=\mathbf{o}_t,\quad
      \widetilde{\mathbf{y}}_t=\widetilde{\mathbf{o}}_t,\\
    \text{Res. space:}\quad
    &\mathbf{m}_s=W_O\mathbf{v}_s,\quad
      \mathbf{r}_t=\hz{}{t},\\
    &\mathbf{y}_t=\Delta_t^{\Attn},\quad
      \widetilde{\mathbf{y}}_t=\widetilde{\Delta}_t^{\Attn}.
  \end{aligned}
  \label{eq:component-scaling-choices}
\end{equation}
For multi-head attention, the value-space quantities are defined separately
for each head before $W_O$. In the residual-space equations,
$\mathbf{v}_s$ denotes the concatenated per-head value message and $W_O$ maps
that concatenation into the residual stream. Consequently, the residual-space
diagonal coefficient is layer-level rather than attributable to one head after
head mixing by $W_O$.
The edited output is then
\begin{equation}
  \widetilde{\mathbf{y}}_t(s^{(\parallel)},s^{(\perp)})
    = s^{(\parallel)}\mathbf{y}_{t,\parallel}
      + s^{(\perp)}\mathbf{y}_{t,\perp}.
  \label{eq:unified-component-scaling}
\end{equation}
The unedited baseline is
$s^{(\parallel)}=s^{(\perp)}=1$; parallel removal sets
$s^{(\parallel)}=0,s^{(\perp)}=1$.
For exclude-self value-space scaling, the same rescaling is applied to
$\mathbf{y}_t=\mathbf{c}_t$ and the fixed $\mathbf{d}_t$ is added afterward,
as in Eq.~\ref{eq:exclude-self-scaling}.

A useful consequence is that a parallel-only edit can be implemented as a
change to the attention diagonal. When $s^{(\perp)}=1$, the edit only changes
the coefficient of the component along $\mathbf{r}_t$. From
Eq.~\ref{eq:unified-component-decomp}, this coefficient is
\[
  c_t=\frac{\mathbf{y}_t^\top\mathbf{r}_t}{\|\mathbf{r}_t\|^2},
  \qquad
  \mathbf{y}_{t,\parallel}=c_t\mathbf{r}_t .
\]
Scaling the parallel component changes this coefficient by
$(s^{(\parallel)}-1)c_t$. If we keep off-diagonal weights fixed and only edit
the diagonal attention entry, then
$\widetilde{\mathbf{y}}_t=\mathbf{y}_t+\Delta A_{tt}\mathbf{m}_t$. Matching
the resulting coefficient along $\mathbf{r}_t$ gives
\begin{equation}
  \Delta A_{tt}
    = (s^{(\parallel)}-1)
      \frac{\mathbf{y}_t^\top\mathbf{r}_t}{\mathbf{m}_t^\top\mathbf{r}_t}.
  \label{eq:generic-diagonal-closed-form}
\end{equation}
Substituting the two choices in Eq.~\ref{eq:component-scaling-choices} into
Eq.~\ref{eq:generic-diagonal-closed-form} gives
\begin{equation}
  \begin{aligned}
  \Delta A_{tt}^{\mathrm{val,raw}}
    &= (s^{(\parallel)}-1)
       \frac{(\sum_{s\le t} A_{ts}\mathbf{v}_s)^\top\mathbf{v}_t}
            {\|\mathbf{v}_t\|^2},\\
  \Delta A_{tt}^{\mathrm{val,excl}}
    &= (s^{(\parallel)}-1)
       \frac{(\sum_{s<t} A_{ts}\mathbf{v}_s)^\top\mathbf{v}_t}
            {\|\mathbf{v}_t\|^2},\\
  \Delta A_{tt}^{\mathrm{res}}
    &= (s^{(\parallel)}-1)
       \frac{(\sum_{s\le t} A_{ts}W_O\mathbf{v}_s)^\top\hz{}{t}}
            {(W_O\mathbf{v}_t)^\top\hz{}{t}}.
  \end{aligned}
  \label{eq:space-diagonal-closed-form}
\end{equation}
Thus parallel scaling in either space can be written as an effective diagonal
change while keeping the off-diagonal messages fixed. For the exclude-self
variant, this diagonal correction is determined only by the non-self aggregate;
the original $A_{tt}\mathbf{v}_t$ term remains present, and no row
renormalization is performed. The residual-space formula solves a different
scalar constraint because its diagonal message $W_O\mathbf{v}_t$ is generally
not the residual reference direction $\hz{}{t}$;
Appendix~\ref{app:method} discusses this distinction.
